\section{Data Augmentation and Feature Engineering}
The task was to improve the detector's robustness against real-world settings,
to still meet the $0.2$ FPR limit.
We decided to continue from the Task~2 checkpoint and tested the model under
altered conditions. This resulted in an FPR of $0.262$, well above the allowed
threshold.

\textbf{Approach}\\
To fix this, we took the transformations mentioned in the task
(scaling, compression, blur) and added a few more common
ones we felt were missing (brightness/contrast jitter, additive noise,
horizontal flips), applying them randomly to each batch during training.
We picked the best checkpoint based on a combined clean+augmented validation
set, so we wouldn't accidentally pick an early epoch that already performs badly
on clean data. We also recalibrated the threshold on combined clean+augmented
calibration data, using a slightly more conservative target
($\tau = 0.14$ instead of $0.15$) to leave extra margin given the noisy FPR
estimate from few real samples.

\textbf{Results}\\
Table~\ref{tab:t2} (last row) shows the augmented model lowers
recall on \texttt{validation\_augmented} from $\CnnAugRecall$ (clean CNN) to
$\AugAugRecall$ at FPR~$\AugAugFpr$, so we are still clearing the $0.6$ target 
and now meet the FPR limit of $0.2$.
The validation performance for the untransformed data is also  
affected, with a recall of 0.756 (previous 0.802) at FPR 0.181 but still
meets all requirements.
